\subsection{P-18. An arbitrary-length trilinear class and a two-boundary cubic gate}
\label{P-18}

\paragraph{Statement.}
Fix distinct nodes $a_1,\ldots,a_n$ over any field $E$.
Suppose that a nonzero multilinear polynomial $f$ has width at
most three in some variable order and is separately homogeneous
of degree one in each block of a partition of all its variables
into three disjoint nonempty consecutive blocks of that same
order. The blocks and order need not be supplied to
the evaluator. Then balanced rational two-pair evaluation
\[
 R_i=t+\frac{z_1}{a_i-y_1}+\frac{z_2}{a_i-y_2}
 \tag{P18.1}
\]
hits $f$. This gives a uniform finite hitting set of at most
\[
 20(n+1)^2
 \tag{P18.2}
\]
points over a working field of size at least
$\max\{2n+1,4\}$, with the explicitly stated extension-input
convention of I-40. The assertion holds in every characteristic.

There is also a precise reduction of the unrestricted cubic
question. If a homogeneous multilinear cubic of width at most
three annihilates (P18.1), choose such an annihilator with the
fewest essential variables and a witnessing order. Its support
triples have their first, second, and third positions in three
consecutive intervals that overlap in at most one variable at
each interface. At least one overlap must occur. Thus at most
two shared boundary variables are the remaining obstruction
in this reduced cubic problem; this statement does not resolve
them or the unrestricted cubic theorem.

\paragraph{Inputs and notation.}
Write $\mathcal I_k$ for the formal vanishing ideal of balanced
rational functions with numerator and denominator degrees at
most $k$. These ideals are prime kernels of rational parameter
substitutions. We use the established HMM degree and exterior
characterizations (P-B20, Definition~9.5, p.~55, and
Theorem~9.13, pp.~61--62), the pole extraction proved in I-37,
and the three-consecutive-degree-block lemma of P-17, part~1.
In particular $\mathcal I_2$ has no nonzero multilinear component
of degree less than three, $\mathcal I_1$ contains no nonzero
linear form, and every partial derivative of a cubic in
$\mathcal I_2$ belongs to $\mathcal I_1$.
The argument below uses those inputs rather than reproving
rational-function ideal membership.

\paragraph{1. A balanced cubic that is trilinear in disjoint blocks cannot have width three.}
Suppose for contradiction that $f\in\mathcal I_2$ satisfies
the stated disjoint-block promise, with blocks $U,V,W$ in
that order. Remove unused variables; this preserves both the
width bound and formal generator membership at the retained
nodes. A nonzero cubic in $\mathcal I_2$ is invariant under
common translation: every positive-power coefficient of
$f(X+s\mathbf1)$ still belongs to $\mathcal I_2$ and has
multilinear degree at most two, so it is zero.

The coefficient of $s$ is the sum of the derivatives in the
three blocks. Those three summands have distinct block degrees
$(0,1,1)$, $(1,0,1)$, and $(1,1,0)$, respectively. Each
summand is therefore zero. Since the degree in each entire
block is one, this also gives separate common-translation
invariance within each block. This uses no division by the
characteristic.

Let $r$ be the coefficient rank across $U\mid(V,W)$. Write
\[
 f=\sum_{i\in U}X_i B_i(V,W),\qquad
 r=\dim\operatorname{span}\{B_i:i\in U\}.
 \tag{P18.3}
\]
Separate translation in $U$ gives $\sum_i B_i=0$. Let $i_*$
be the last variable of $U$ in the witnessing order. It is
essential, so $B_{i_*}\ne0$, and the earlier $B_i$ span the
full $r$-dimensional space because their sum is $-B_{i_*}$.
At the cut just before $X_{i_*}$, the rows indexed by earlier
single variables span these $r$ dimensions. The constant-prefix
row is $X_{i_*}B_{i_*}$, independent of them by its degree in
$X_{i_*}$. Hence this cut has rank at least $r+1$, and
\[
 r\le2.
 \tag{P18.4}
\]
If $U$ had only one essential variable, the same sum relation
would already make $f=0$. Applying the identical argument in
the reversed order gives rank at most two across
$(U,V)\mid W$.

Neither rank can be one. A rank-one extreme block cut would
factor $f$ into a nonzero linear form and a nonzero quadratic
in disjoint variables. Primality of $\mathcal I_2$ would put
one factor in that ideal, contrary to its minimum multilinear
degree. Thus both ranks are exactly two. Choose independent
linear-form vectors $L(U)$ and $N(W)$, each of length two.
The two coefficient-space factorizations give
\[
 f=L(U)^{\mathsf T}M(V)N(W),\qquad
 M(V)=\sum_{j\in V}X_j M_j,\quad M_j\in E^{2\times2}.
 \tag{P18.5}
\]
All middle variables are essential, so $M_j\ne0$ for each.
Pole extraction gives
$\partial_j f=L^{\mathsf T}M_jN\in\mathcal I_1$.
If $M_j$ had rank one, this nonzero quadratic would be a
product of two nonzero linear forms, which a prime ideal
without nonzero linear forms cannot contain. Therefore every
$M_j$ is invertible. Separate translation in $V$ gives
$\sum_jM_j=0$, so there are at least two middle variables.

Cut between two consecutive middle variables, leaving a
nonempty middle prefix and suffix. The part using a prefix
variable has left degree two and right degree one, and its
coefficient rank is two: one invertible $M_j$ in the prefix,
together with the independent $L,N$, already attains that
rank. The part using a suffix variable has left degree one
and right degree two and likewise has rank two. Their degree
blocks have disjoint row and column indices. The total cut
rank is at least four, a contradiction. This proves the
nonvanishing assertion for all $n$ and every characteristic.

The denominator-safe interpolation argument of I-40, part~6,
now applies at total degree $D=3$. Each selector degree is at
most $n$ after multiplication by
$\prod_i(a_i-y_1)(a_i-y_2)$, and the three-amplitude total
degree is at most three. Its grid size is
$(n+1)^2\binom{6}{3}=20(n+1)^2$, with the field size stated
in (P18.2). The grid does not require the witnessing blocks
or variable order.

\paragraph{2. A smallest unrestricted counterexample has at most two shared boundary variables.}
Suppose that a cubic counterexample exists, choose one with
the fewest essential variables, and fix a width-three order.
For any proper prefix cut $A\sqcup B$, let $C_d$ be its
coefficient block of left degree $d$, $0\le d\le3$.
We first note a direct consequence of the established exterior
characterization. Decode coefficients by
\[
 W=\sum_{|S|=3}\frac{c_S}{\Delta_S}e_S,
 \qquad
 W\in\bigwedge^3\ker[\mathbf1;\mathbf a;\mathbf a^{\circ2}].
 \tag{P18.6}
\]
The three contraction equations connect only adjacent left
exterior degrees: an output two-vector of left degree $e$
receives a left contraction from degree $e+1$ or a right
contraction from degree $e$. Therefore a zero degree block
separates the kernel equations on its two sides. In particular,
if $C_1=0$, the degree-zero part alone still satisfies all three
contractions, so the polynomial represented by $C_0$ belongs
to $\mathcal I_2$. If $C_2=0$, the analogous statement holds
for $C_3$. This is the adjacent-block separation mechanism
of HMM Proposition~8.6, printed pp.~44--45, applied here to
the balanced cubic's stated contraction equations.

A nonzero $C_0$ or $C_3$ so separated is obtained by setting
one side of the variables to zero. Its width cannot increase,
and it has fewer essential variables, contradicting the
choice of the counterexample. Consequently $C_0\ne0$ forces
$C_1\ne0$. P-17's three-block lemma then forces $C_2=0$;
minimality forces $C_3=0$. Symmetrically $C_3\ne0$ allows
only $C_2,C_3$. If both extremes vanish, only $C_1,C_2$ remain.
For every prefix, all occupied degree blocks are thus contained
in two adjacent positions.

Equivalently, any two support triples have prefix counts that
differ by at most one. Define
\[
 \begin{aligned}
 \alpha&=\max\{\text{first position of a support triple}\},\\
 \beta&=\min\{\text{second position of a support triple}\},\\
 \gamma&=\max\{\text{second position of a support triple}\},\\
 \delta&=\min\{\text{third position of a support triple}\}.
 \end{aligned}
 \tag{P18.7}
\]
Then $\alpha\le\beta$: otherwise the prefix ending at $\beta$
contains two variables of one support triple and none of a
triple whose first position is $\alpha$. Similarly
$\gamma\le\delta$, since the prefix ending at $\delta$
would otherwise contain three variables of one triple and
at most one of a triple whose second position is $\gamma$.
Every monomial therefore uses one first position in
$[1,\alpha]$, one second position in $[\beta,\gamma]$, and
one third position in $[\delta,n]$.

No variable can lie in a gap between these intervals, since
all variables are essential. The intervals are consecutive,
and the only possible intersections are the singleton
$\{\alpha\}$ if $\alpha=\beta$ and the singleton
$\{\gamma\}$ if $\gamma=\delta$. If neither intersection
occurs, the polynomial is trilinear in three disjoint
consecutive blocks and part~1 excludes it. Thus any smallest
counterexample has one or two shared boundary variables.
This proves the reduction.\hfill$\square$

\paragraph{Scope and the unresolved boundary cases.}
The trilinear promise requires three consecutive blocks in
one witnessing width-three order. It does not assert that an
arbitrary linear change of coordinates preserves such an order
or the coefficient-cut ranks. The bound is a concrete
arbitrary-length restricted sampler, not polynomial-in-growing-width
PIT. P-19 subsequently proves the full cubic two-pair theorem,
using the interval reduction established here.

Setting a shared boundary variable to zero cannot simply be
used to remove the overlap: it need not preserve the balanced
ideal. Imposing a zero at one node lowers the rational numerator
budget after the corresponding diagonal rescaling. The coupled
quadratic partials at the one or two boundary variables are
the remaining gate at this stage of the argument; P-19
subsequently closes both boundary cases. The separate audit
\texttt{cubic\_ordered\_gate.md} also records a direct
all-characteristic cubic prolongation consequence of the HMM
kernels; it is not needed for part~1. No historical priority
claim is made.
